% ============================================================================
% 01-metodologia.tex — Methods (module 2/5) com fluxograma da validação cruzada
% ============================================================================
\section{Methods}
\label{sec:methods}

\subsection{System under test}
\label{sec:system}
All experiments ran inside the OpenCode Ecosystem Core, a multi-agent
ecosystem with 206 catalog agents, an A2A Blackboard, a metacognitive bus
(MetaBus), a Trust Engine and a strict SDD/TDD policy. The orchestrator
(MarceloClaro) executes the cycle Perceive$\rightarrow$Specify$\rightarrow$
Delegate$\rightarrow$Execute$\rightarrow$Verify$\rightarrow$Reflect. External
LLMs are invoked through the \texttt{opencode run} CLI as real subprocesses;
no simulation or mocked completion is used (Section~\ref{sec:anti-overclaim}).

\subsection{Corpus}
\label{sec:corpus}
We use a canonical IMO-AnswerBench sample (R442, 4 problems) plus five
IMO-style auxiliary problems with deterministically verifiable answers
(R503), totalling $n=9$: algebra-001 (quotient floor, answer 3),
algebra-004 (inequality, $2^{u-2}$), number-theory-001 ($p^3-q^5=(p+q)^2$,
answer $(7,3)$), combinatorics-001 (one local maximum, $2^{n-1}$),
algebra-002 ($2^n>n^2$, $n>1$, answer 5), number-theory-002 ($p^2 \mid 2^p+1$,
$\{3\}$), number-theory-003 ($v_3(2023!)=1006$, Legendre), combinatorics-002
(two local maxima in $S_5$, verified 88 by independent enumeration),
geometry-001 (2023-gon diagonals, 2043230). The auxiliary items are
IMO-style derivatives, not official shortlist entries; we label them
accordingly to avoid attribution errors.

\subsection{Anti-leak deterministic solver}
\label{sec:solver}
The original IMO harness \citep{imo_harness_local} echoed the expected
short answer and self-awarded grade 7 --- a leak. We implemented
\texttt{RealIMOSolver}, a fully deterministic solver (enumeration, Legendre
formula, polygon formula) whose output never contains the expected answer
before its ``Metodo:'' marker, and verified independence by round-trip tests
(none of the 9 solutions leaks the answer). The solver's own cross-check
scores 8/9 on the augmented corpus (algebra-004 is numerical-only,
``partial'', never grade 7).

\subsection{Conditions}
\label{sec:conditions}
Three conditions on the same NT-001 probe problem:
C0 raw prompt; C1 context-only; C2 16-stage MASWOS scaffold (reading,
hypotheses, verification, strict output format). Composite score
$S = 0.40A + 0.30V + 0.20C + 0.10D$, where $A$=accuracy, $V$=normalized
latency $1-(\ell-20)/90$ (clamped to $[0,1]$), $C$=format conformity,
$D$=availability.

\subsection{Cross-validation protocol (R503)}
\label{sec:crossval}
After the R499 single-problem ranking, we implemented automatic routing
(R500) and then re-ran the experiment after a full host restart (R503).
The cross-validated protocol is a 9-problem $\times$ 3-model full factorial
design (no held-out split; all models see all problems, exactly once, in C2).
We report:
\begin{itemize}
  \item per-model accuracy $k/n$ with exact Clopper--Pearson 95\% confidence
  intervals (binomial);
  \item paired McNemar tests with continuity correction between model pairs;
  \item Cohen's $\kappa$ (agreement beyond chance) for the two most
  competent models;
  \item per-model latency distributions (ok vs.\ timeout).
\end{itemize}
Figure~\ref{fig:flow} shows the cross-validation flow executed per model.

% ============================================================================
% fig-flow-valcross.tex — Fluxograma da validação cruzada 9×3 (Figura 4)
% ============================================================================
\begin{figure}[htbp]
\centering
\begin{tikzpicture}[
  node distance=0.8cm and 1.2cm,
  caixa/.style={rectangle, draw=cororo, thick, rounded corners=2pt,
                fill=cororo!8, minimum width=3.0cm, minimum height=0.9cm,
                align=center, font=\footnotesize},
  decisao/.style={diamond, draw=coramarelo, thick, fill=coramarelo!10,
                  aspect=2, align=center, font=\footnotesize},
  seta/.style={-{Stealth[length=3mm]}, thick, color=corcinza}
]
  \node[caixa] (corpus) {Corpus IMO\\ 9 problemas (\S2.2)};
  \node[decisao, right=of corpus] (modelo) {modelo\\ mimo / big-pickle /\\ nemotron};
  \node[caixa, below=of modelo] (c2) {Condição C2\\ scaffold MASWOS (16 estágios)};
  \node[caixa, below=of c2] (cli) {CLI real\\ \texttt{opencode run}};
  \node[decisao, below=of cli] (fmt) {saída\\ conforme?};
  \node[caixa, below=of fmt] (ext) {Extração de resposta\\ (regex por problema)};
  \node[caixa, below=of ext] (estat) {Estatística\\ acc + IC Clopper-Pearson\\ McNemar pareado\\ Cohen's $\kappa$ \\ latência};

  \draw[seta] (corpus) -- (modelo);
  \draw[seta] (modelo) -- (c2);
  \draw[seta] (c2) -- (cli);
  \draw[seta] (cli) -- (fmt);
  \draw[seta] (fmt) -- node[right, font=\tiny]{sim} (ext);
  \draw[seta, color=corvermelho] (fmt) -- ++(2.6,0) node[right, font=\tiny, align=left, color=corvermelho]{não: vazio/timeout\\ $\rightarrow$ incorreto};
  \draw[seta] (ext) -- (estat);
  \draw[seta, dashed] (estat) -- ++(-3.2,-1.2) node[left, font=\tiny]{repete p/ cada modelo};
\end{tikzpicture}
\caption{Fluxograma da validação cruzada 9 $\times$ 3 (R503): cada modelo
percorre os 9 problemas na condição C2; respostas são extraídas por regex por
problema e pontuadas; estatísticas de confiança/incerteza são calculadas
sobre a matriz completa.}
\label{fig:flow}
\end{figure}

\subsection{Cost budget}
Each LLM call is a real CLI subprocess with a model-specific timeout
(mimo 100~s, big-pickle 110~s, nemotron 150~s). All raw outputs, timings and
statuses are stored incrementally (JSON) inside the repository for audit.